\documentclass[12pt,a4paper]{report}
\usepackage[utf8]{inputenc}
\usepackage[T1]{fontenc}
\usepackage[french]{babel}
\usepackage{lmodern}
\usepackage[margin=2.5cm]{geometry}
\usepackage{amsmath,amssymb}
\usepackage{booktabs}
\usepackage{array}
\usepackage{enumitem}
\usepackage{xcolor}
\usepackage{colortbl}
\usepackage{listings}
\usepackage{tikz}
\usetikzlibrary{positioning}
\usepackage{pgfplots}
\pgfplotsset{compat=1.16}
\usepackage{caption}
\usepackage{float}
\usepackage{etoolbox}
\usepackage[colorlinks=true,linkcolor=blue!70!black,urlcolor=blue!70!black]{hyperref}

\hypersetup{pdftitle={Benchmark arbre rouge-noir}, pdfauthor={BENNOUR RAFAA}}
\captionsetup{labelsep=endash, font=small}
\captionsetup[table]{position=above}
\emergencystretch=2em
\widowpenalty=10000 \clubpenalty=10000
\setlist{itemsep=2pt}

% ---------------- Code ----------------
\definecolor{fondcode}{HTML}{F6F6F3}
\definecolor{bordcode}{HTML}{CCCCCC}
\definecolor{motcle}{HTML}{1F5FA8}
\definecolor{commentaire}{HTML}{6E6E6E}
\definecolor{ok}{HTML}{D4EDD4}
\definecolor{ko}{HTML}{F6D1CC}
\definecolor{cmain}{HTML}{1F5FA8}
\definecolor{cpq}{HTML}{D35400}
\lstset{
  language=C++, basicstyle=\ttfamily\footnotesize,
  keywordstyle=\color{motcle}\bfseries, commentstyle=\color{commentaire}\itshape,
  backgroundcolor=\color{fondcode}, frame=single, rulecolor=\color{bordcode},
  columns=fullflexible, keepspaces=true, showstringspaces=false,
  xleftmargin=4pt, xrightmargin=4pt, aboveskip=8pt, belowskip=8pt,
  literate={é}{{\'e}}1 {è}{{\`e}}1 {à}{{\`a}}1 {ê}{{\^e}}1 {ç}{{\c{c}}}1
           {ù}{{\`u}}1 {É}{{\'E}}1 {î}{{\^i}}1 {ô}{{\^o}}1 {û}{{\^u}}1 {â}{{\^a}}1
}
\lstdefinestyle{sortie}{language={}, basicstyle=\ttfamily\scriptsize, keywordstyle={}, backgroundcolor=\color{white}}
\BeforeBeginEnvironment{lstlisting}{\par\noindent\begin{minipage}{\linewidth}}
\AfterEndEnvironment{lstlisting}{\end{minipage}\par}

% ---------------- Arbres ----------------
\tikzset{
  noeud/.style={circle, draw=black!70, fill=blue!8, minimum size=7mm, inner sep=0pt, font=\small},
  racine/.style={noeud, fill=orange!30},
  level distance=11mm,
}
\tikzset{
  rn/.style={circle, draw=black, fill=black!85, text=white, minimum size=7mm, inner sep=0pt, font=\small},
  rr/.style={circle, draw=red!70!black, fill=red!75!black, text=white, minimum size=7mm, inner sep=0pt, font=\small},
  nil/.style={rectangle, draw=black!50, fill=black!10, minimum size=4mm, inner sep=1pt, font=\tiny},
}
\definecolor{cabr}{HTML}{7F7F7F}
\definecolor{crn}{HTML}{C0392B}
\definecolor{cmap}{HTML}{1F5FA8}
\definecolor{cum}{HTML}{1E8449}

% =====================================================================
% =====================================================================
\begin{document}

\begin{titlepage}
  \centering
  \vspace*{2cm}
  {\huge Benchmark Comparatif\par}
  \vspace{0.4cm}
  {\LARGE Arbre rouge-noir\par}
  \vspace{0.3cm}
  {\Large ABR, rouge-noir, \texttt{std::map} et \texttt{std::unordered\_map}\par}
  \vspace{1.3cm}
  {\large Insertion, recherche et suppression de clés entières :\\[2pt]
  analyse de performance comparative\par}
  \vspace{1.3cm}
  \begin{tikzpicture}[level distance=12mm, level 1/.style={sibling distance=34mm}, level 2/.style={sibling distance=17mm}]
    \node[rn] {48}
      child { node[rn] {7}
        child { node[rr] {2} }
        child { node[rr] {17} } }
      child { node[rn] {82}
        child { node[rr] {58} }
        child { node[rr] {84} } };
  \end{tikzpicture}
  \vspace{1.3cm}

  \begin{tabular}{r@{\quad}l}
    \textbf{Auteur :} & BENNOUR RAFAA \\
    \textbf{Formation :} & ILISI -- 2\textsuperscript{e} année \\
    \textbf{Date :} & 11 novembre 2026 \\
    \textbf{Langage :} & C++17 \\
    \textbf{Compilateur :} & g++ 15.2 \texttt{-O2} \\
  \end{tabular}
\end{titlepage}

\setcounter{tocdepth}{1}
\tableofcontents

% =====================================================================
\chapter{Introduction}

\section{Contexte et objectifs}
Un arbre binaire de recherche (ABR) range les valeurs plus petites à gauche et les plus grandes à droite. Chercher, insérer ou supprimer une valeur coûte autant d'étapes que la \textbf{hauteur} de l'arbre. Le problème : si les valeurs arrivent triées, l'ABR devient une liste et sa hauteur vaut $n$.

L'\textbf{arbre rouge-noir} corrige ce défaut : il donne une couleur à chaque nœud et se répare après chaque insertion et chaque suppression, par des recolorations et des rotations. Sa hauteur reste toujours inférieure à $2\log_2(n+1)$.

Ce rapport compare quatre structures qui permettent de ranger des clés, de les retrouver et de les supprimer :
\begin{description}[leftmargin=0.5cm, style=nextline]
  \item[ABR] arbre binaire de recherche simple, \textbf{non équilibré}, écrit à la main ;
  \item[Rouge-noir] arbre rouge-noir \textbf{écrit à la main} (insertion, suppression, rotations) ;
  \item[\texttt{std::map}] l'arbre rouge-noir de la bibliothèque standard ;
  \item[\texttt{std::unordered\_map}] la table de hachage de la bibliothèque standard.
\end{description}

\subsection{Objectifs de l'étude}
\begin{enumerate}
  \item Implémenter un arbre rouge-noir complet (insertion, suppression, rotations, recoloration) et vérifier qu'il respecte ses règles.
  \item Mesurer l'effet de l'équilibrage, en insérant les clés dans un ordre \textbf{aléatoire} puis \textbf{croissant}.
  \item Comparer notre arbre rouge-noir avec celui de la bibliothèque standard (\texttt{std::map}).
  \item Situer les arbres face à une table de hachage (\texttt{std::unordered\_map}).
\end{enumerate}

\section{L'arbre rouge-noir}

\subsection{Les cinq règles}
\begin{enumerate}
  \item Chaque nœud est rouge ou noir.
  \item La racine est noire.
  \item Les feuilles vides (NIL, \texttt{nullptr}) sont considérées comme noires.
  \item Un nœud rouge n'a jamais d'enfant rouge.
  \item Depuis un nœud, tous les chemins vers les feuilles passent par le même nombre de nœuds noirs.
\end{enumerate}
Les règles 4 et 5 garantissent l'équilibre : tous les chemins ont le même nombre $b$ de nœuds noirs, et on ne peut pas placer deux rouges à la suite, donc le chemin le plus long fait au plus $2b$ nœuds : il est \textbf{au plus deux fois plus long que le plus court}.

\subsection{L'insertion}
On insère comme dans un ABR, avec un nœud \textbf{rouge}. Si son parent est rouge, la règle 4 est cassée et on répare selon la couleur de l'\textbf{oncle} (le frère du parent) :

\begin{center}
\small
\begin{tabular}{p{0.14\linewidth}p{0.30\linewidth}p{0.46\linewidth}}
  \toprule
  Cas & Situation & Réparation \\
  \midrule
  Cas 1 & oncle rouge & parent et oncle deviennent noirs, grand-parent rouge ; on recommence avec le grand-parent \\
  Cas 2 & oncle noir, nœud en zigzag & rotation sur le parent pour se ramener au cas 3 \\
  Cas 3 & oncle noir, nœud en ligne droite & rotation sur le grand-parent, puis échange des couleurs ; terminé \\
  \bottomrule
\end{tabular}
\end{center}

Une \textbf{rotation} ne modifie que trois liens, elle coûte $O(1)$ :
\begin{center}
\begin{tikzpicture}[level distance=10mm, level 1/.style={sibling distance=16mm}, level 2/.style={sibling distance=10mm}]
  \node[noeud] {x}
    child { node[nil] {a} }
    child { node[noeud] {y}
      child { node[nil] {b} }
      child { node[nil] {c} } };
  \node at (2.6,-1) {\Large $\longrightarrow$};
  \node at (2.6,-0.6) {\small rotation gauche};
  \begin{scope}[xshift=5.2cm]
  \node[noeud] {y}
    child { node[noeud] {x}
      child { node[nil] {a} }
      child { node[nil] {b} } }
    child { node[nil] {c} };
  \end{scope}
\end{tikzpicture}
\end{center}
Une insertion fait au plus deux rotations et $O(\log n)$ recolorations : elle coûte $O(\log n)$.

\subsection{La suppression}
On retire d'abord le nœud comme dans un ABR : s'il a deux enfants, son \textbf{successeur} (le minimum du sous-arbre droit) prend sa place et sa couleur. Si le nœud qui quitte vraiment sa place est \textbf{rouge}, aucune règle n'est cassée. S'il est \textbf{noir}, il manque un noir sur certains chemins (règle 5) : le nœud $x$ qui prend la place reçoit un « noir en trop » (on dit qu'il est \textbf{doublement noir}). Si $x$ est rouge, il devient simplement noir. Sinon, on répare selon le \textbf{frère} de $x$ et ses enfants (les \textbf{neveux}) :

\begin{center}
\small
\begin{tabular}{p{0.14\linewidth}p{0.36\linewidth}p{0.40\linewidth}}
  \toprule
  Cas & Situation & Réparation \\
  \midrule
  Cas 1 & frère rouge & rotation sur le parent et échange des couleurs : le frère devient noir, on passe aux cas 2 à 4 \\
  Cas 2 & frère noir, deux neveux noirs & le frère devient rouge ; le noir en trop remonte au parent \\
  Cas 3 & frère noir, neveu proche rouge, neveu lointain noir & rotation sur le frère pour se ramener au cas 4 \\
  Cas 4 & frère noir, neveu lointain rouge & rotation sur le parent et recoloration ; terminé \\
  \bottomrule
\end{tabular}
\end{center}
Une suppression fait au plus trois rotations et $O(\log n)$ recolorations : elle coûte aussi $O(\log n)$, mais elle fait plus de travail qu'une insertion (chercher le successeur, puis réparer).

\section{Jeux de données}
\begin{itemize}
  \item Les clés sont les entiers $0, 1, \dots, N-1$, sans doublon, avec $N = 10^3$, $10^4$, $2\times10^4$, $10^5$ et $10^6$.
  \item \textbf{Ordre aléatoire} : les clés sont mélangées (graine fixe).
  \item \textbf{Ordre croissant} : les clés sont insérées triées. C'est le pire cas de l'ABR, qui n'est mesuré que jusqu'à $N = 20\,000$ (au-delà, son coût en $O(N^2)$ devient trop long).
  \item Démonstration : insertion de 7 58 48 2 17 84 82, puis suppression de 2, 58, 48, 7 et 99 (absent).
\end{itemize}

\section{Opérations mesurées}
\begin{itemize}
  \item \textbf{Insertion} : insérer les $N$ clés dans une structure vide.
  \item \textbf{Recherche} : chercher les $N$ clés, dans un autre ordre aléatoire. Toutes doivent être trouvées.
  \item \textbf{Suppression} : supprimer les $N$ clés, dans ce même ordre aléatoire, jusqu'à vider la structure. Toutes doivent être supprimées.
  \item \textbf{Hauteur} : nombre de niveaux des deux arbres écrits à la main, mesuré avant les suppressions (elle n'est pas accessible pour \texttt{std::map}).
\end{itemize}
Chaque résultat est un temps moyen par opération, en nanosecondes (ns).

\section{Complexités théoriques}
\begin{table}[H]
\centering
\caption{Complexités (cas moyen, avec le pire cas entre parenthèses).}
\small
\begin{tabular}{lcc}
  \toprule
  Structure & Insertion, recherche, suppression & Ordre des clés \\
  \midrule
  ABR & $O(\log n)$ \quad ($O(n)$) & trié \\
  Rouge-noir / \texttt{map} & $O(\log n)$ garanti & trié \\
  \texttt{unordered\_map} & $O(1)$ \quad ($O(n)$) & aucun \\
  \bottomrule
\end{tabular}
\end{table}

% =====================================================================
\chapter{Méthodologie}

\section{Environnement}
\begin{itemize}
  \item Ordinateur portable : Intel Core i7-8550U, 8~Go de mémoire, Windows 11.
  \item Compilateur MinGW g++ 15.2, options \texttt{-O2 -std=c++17}.
  \item Les temps dépendent de la machine ; les tendances restent les mêmes d'un PC à l'autre.
\end{itemize}

\section{Protocole de mesure}
\begin{itemize}
  \item Chronomètre : \texttt{std::chrono::steady\_clock}.
  \item Chaque mesure est faite 3 fois et l'on garde le meilleur temps.
  \item Les quatre structures reçoivent exactement les mêmes clés et les mêmes requêtes.
  \item \textbf{Vérifications} avant les mesures, sur $100\,000$ clés : l'arbre rouge-noir respecte ses règles après les insertions (aléatoires et croissantes), puis après la suppression d'une clé sur deux ; les bonnes clés sont encore présentes ; et l'arbre est vide une fois toutes les clés supprimées.
  \item \textbf{Vérifications} pendant les mesures : chaque clé cherchée est trouvée et chaque clé supprimée était bien présente.
\end{itemize}

\section{Implémentation}

\subsection{Le nœud}
Le nœud rouge-noir a un pointeur vers son \textbf{parent}, nécessaire pour remonter vers le grand-parent, l'oncle ou le frère.
\begin{lstlisting}
enum Couleur { ROUGE, NOIR };
struct NoeudRN {
    int valeur;
    Couleur couleur = ROUGE;               // tout nouveau noeud est rouge
    NoeudRN *gauche = nullptr, *droite = nullptr, *parent = nullptr;
};
static Couleur couleurDe(NoeudRN* n) { return n == nullptr ? NOIR : n->couleur; }
\end{lstlisting}

\subsection{La rotation gauche}
\begin{lstlisting}
void rotationGauche(NoeudRN* x) {
    NoeudRN* y = x->droite;
    x->droite = y->gauche;                       // 1. b passe à droite de x
    if (y->gauche != nullptr) y->gauche->parent = x;
    y->parent = x->parent;                       // 2. y prend la place de x
    if (x->parent == nullptr)        racine = y;
    else if (x == x->parent->gauche) x->parent->gauche = y;
    else                             x->parent->droite = y;
    y->gauche = x;                               // 3. x passe à gauche de y
    x->parent = y;
}
\end{lstlisting}

\subsection{La réparation après insertion}
Le code ci-dessous traite le cas où le parent est l'enfant gauche du grand-parent ; l'autre côté est son miroir (gauche et droite inversés).
\begin{lstlisting}
while (z != racine && couleurDe(z->parent) == ROUGE) {
    NoeudRN* p = z->parent;  NoeudRN* g = p->parent;
    if (p == g->gauche) {
        NoeudRN* oncle = g->droite;
        if (couleurDe(oncle) == ROUGE) {           // CAS 1 : recolorer
            p->couleur = NOIR; oncle->couleur = NOIR; g->couleur = ROUGE;
            z = g;                                 // on remonte
        } else {
            if (z == p->droite) {                  // CAS 2 : zigzag
                z = p; rotationGauche(z); p = z->parent;
            }
            p->couleur = NOIR; g->couleur = ROUGE; // CAS 3 : ligne droite
            rotationDroite(g);
        }
    } else { /* miroir */ }
}
racine->couleur = NOIR;
\end{lstlisting}

\subsection{La suppression}
Les nœuds sont déplacés en changeant les pointeurs (\texttt{transplanter}), jamais en recopiant les valeurs : seul le nœud supprimé disparaît.
\begin{lstlisting}
Couleur couleurRetiree = z->couleur;
if (z->gauche == nullptr)       { x = z->droite; parentX = z->parent; transplanter(z, z->droite); }
else if (z->droite == nullptr)  { x = z->gauche; parentX = z->parent; transplanter(z, z->gauche); }
else {                                       // deux enfants : le successeur y prend la place
    NoeudRN* y = minimum(z->droite);
    couleurRetiree = y->couleur;  x = y->droite;
    if (y->parent == z) parentX = y;
    else { parentX = y->parent; transplanter(y, y->droite);
           y->droite = z->droite; y->droite->parent = y; }
    transplanter(z, y);
    y->gauche = z->gauche; y->gauche->parent = y;
    y->couleur = z->couleur;
}
delete z;
if (couleurRetiree == NOIR) reparerSuppression(x, parentX);  // il manque un noir
\end{lstlisting}

\subsection{La réparation après suppression}
Comme pour l'insertion, seul le côté gauche est montré ; l'autre côté est son miroir. Le nœud $x$ peut être vide, c'est pourquoi on garde aussi son parent $p$.
\begin{lstlisting}
while (x != racine && couleurDe(x) == NOIR) {
    if (x == p->gauche) {
        NoeudRN* frere = p->droite;
        if (couleurDe(frere) == ROUGE) {                      // CAS 1
            frere->couleur = NOIR; p->couleur = ROUGE;
            rotationGauche(p); frere = p->droite;
        }
        if (couleurDe(frere->gauche) == NOIR && couleurDe(frere->droite) == NOIR) {
            frere->couleur = ROUGE; x = p; p = x->parent;     // CAS 2 : on remonte
        } else {
            if (couleurDe(frere->droite) == NOIR) {           // CAS 3
                frere->gauche->couleur = NOIR; frere->couleur = ROUGE;
                rotationDroite(frere); frere = p->droite;
            }
            frere->couleur = p->couleur; p->couleur = NOIR;   // CAS 4 : terminé
            frere->droite->couleur = NOIR;
            rotationGauche(p); x = racine;
        }
    } else { /* miroir */ }
}
if (x != nullptr) x->couleur = NOIR;
\end{lstlisting}

\subsection{Organisation du projet}
\begin{center}
\small
\begin{tabular}{ll}
  \toprule
  Fichier & Rôle \\
  \midrule
  \texttt{ArbreBinaireRecherche.h} & ABR non équilibré \\
  \texttt{ArbreRougeNoir.h} & arbre rouge-noir et vérification des règles \\
  \texttt{DemoRBT.cpp} & insertion puis suppression pas à pas \\
  \texttt{Benchmark.cpp} & comparaison des quatre structures \\
  \bottomrule
\end{tabular}
\end{center}

% =====================================================================
\chapter{Résultats : démonstration}

\section{Insertion}
Pour la suite 7 58 48 2 17 84 82, le programme affiche l'arbre ci-dessous (racine à gauche, R = rouge, N = noir). C'est le même arbre que celui obtenu à la main.
\begin{lstlisting}[style=sortie]
            84R
      82N
            58R
48N
            17R
      7N
            2R
Regles respectees : oui
\end{lstlisting}

\begin{figure}[H]
\centering
\begin{tikzpicture}[level distance=12mm, level 1/.style={sibling distance=36mm}, level 2/.style={sibling distance=18mm}]
  \node[rn] {48}
    child { node[rn] {7} child { node[rr] {2} } child { node[rr] {17} } }
    child { node[rn] {82} child { node[rr] {58} } child { node[rr] {84} } };
\end{tikzpicture}
\caption{Arbre après les insertions : chaque chemin compte deux nœuds noirs.}
\end{figure}
Les étapes : 48 provoque un zigzag (cas 2 puis cas 3), 2 une recoloration (cas 1), 17 et 84 ne demandent rien, et 82 un nouveau zigzag.

\section{Suppression}
\begin{center}
\small
\begin{tabular}{lp{0.70\linewidth}}
  \toprule
  Suppression & Ce qui se passe \\
  \midrule
  2 et 58 & feuilles rouges : on les retire, aucune règle n'est cassée \\
  48 (racine) & deux enfants : le successeur 82 prend sa place et sa couleur noire ; 82 était noir, donc il manque un noir sous lui, mais 84, qui monte à sa place, est rouge : il devient noir \\
  7 & un seul enfant, 17 (rouge), qui monte et devient noir \\
  99 & absente : la fonction renvoie \texttt{false}, l'arbre ne change pas \\
  \bottomrule
\end{tabular}
\end{center}

\begin{figure}[H]
\centering
\begin{tikzpicture}[level distance=12mm, level 1/.style={sibling distance=26mm}]
  \node[rn] {82}
    child { node[rn] {17} }
    child { node[rn] {84} };
\end{tikzpicture}
\caption{Arbre final après les suppressions : les règles sont toujours respectées.}
\end{figure}
Cette démonstration ne passe que par les cas simples. Les quatre cas de réparation (frère rouge, neveux noirs, etc.) sont exercés par la vérification du benchmark, qui supprime $50\,000$ clés dans un ordre aléatoire et contrôle ensuite les règles et le contenu de l'arbre.

% =====================================================================
\chapter{Résultats : clés en ordre aléatoire}

\section{Données brutes}
\begin{table}[H]
\centering
\caption{Ordre aléatoire : temps moyen par opération (ns). En vert, l'arbre le plus rapide.}
\newcommand{\g}{\cellcolor{ok}}
\small
\begin{tabular}{rlrrrr}
  \toprule
  $N$ & Structure & Hauteur & Insertion & Recherche & Suppression \\
  \midrule
  1\,000 & ABR & 22 & 429.3 & 70.4 & 287.3 \\
         & Rouge-noir & 12 & 414.5 & \g 50.5 & 328.3 \\
         & \texttt{map} & -- & \g 302.1 & 107.6 & \g 225.8 \\
         & \texttt{unordered\_map} & -- & 418.9 & 17.6 & 111.6 \\
  \midrule
  10\,000 & ABR & 29 & 509.8 & 125.4 & \g 250.4 \\
          & Rouge-noir & 16 & 434.6 & \g 94.4 & 292.5 \\
          & \texttt{map} & -- & \g 359.1 & 184.2 & 343.6 \\
          & \texttt{unordered\_map} & -- & 233.4 & 18.0 & 134.2 \\
  \midrule
  20\,000 & ABR & 33 & \g 411.4 & 444.9 & 494.5 \\
          & Rouge-noir & 18 & 445.8 & \g 104.5 & \g 418.6 \\
          & \texttt{map} & -- & 490.3 & 242.1 & 766.8 \\
          & \texttt{unordered\_map} & -- & 306.7 & 38.2 & 142.9 \\
  \midrule
  100\,000 & ABR & 41 & \g 550.0 & 303.9 & \g 499.3 \\
           & Rouge-noir & 20 & 628.2 & \g 242.1 & 695.1 \\
           & \texttt{map} & -- & 669.5 & 601.9 & 689.0 \\
           & \texttt{unordered\_map} & -- & 346.4 & 34.6 & 262.1 \\
  \midrule
  1\,000\,000 & ABR & 50 & 1\,551.0 & 895.5 & \g 1\,401.5 \\
              & Rouge-noir & 24 & \g 1\,365.1 & \g 708.1 & 1\,612.8 \\
              & \texttt{map} & -- & 1\,397.9 & 1\,386.4 & 1\,822.6 \\
              & \texttt{unordered\_map} & -- & 620.4 & 129.2 & 708.8 \\
  \bottomrule
\end{tabular}
\end{table}

\section{Visualisation}
\begin{figure}[H]
\centering
\begin{tikzpicture}
\begin{loglogaxis}[width=0.68\linewidth, height=6.5cm, xlabel={$N$}, ylabel={ns par recherche},
  legend style={font=\small, at={(1.03,0.5)}, anchor=west}, grid=major, grid style={black!12}]
  \addplot[cabr, thick, mark=triangle*] coordinates {(1000,70.4) (10000,125.4) (20000,444.9) (100000,303.9) (1000000,895.5)};
  \addplot[crn, thick, mark=*] coordinates {(1000,50.5) (10000,94.4) (20000,104.5) (100000,242.1) (1000000,708.1)};
  \addplot[cmap, thick, mark=square*] coordinates {(1000,107.6) (10000,184.2) (20000,242.1) (100000,601.9) (1000000,1386.4)};
  \addplot[cum, thick, mark=diamond*] coordinates {(1000,17.6) (10000,18.0) (20000,38.2) (100000,34.6) (1000000,129.2)};
  \legend{ABR, Rouge-noir, \texttt{map}, \texttt{unordered\_map}}
\end{loglogaxis}
\end{tikzpicture}
\caption{Recherche, clés en ordre aléatoire (échelles logarithmiques).}
\end{figure}

\begin{figure}[H]
\centering
\begin{tikzpicture}
\begin{loglogaxis}[width=0.68\linewidth, height=6.5cm, xlabel={$N$}, ylabel={ns par suppression},
  legend style={font=\small, at={(1.03,0.5)}, anchor=west}, grid=major, grid style={black!12}]
  \addplot[cabr, thick, mark=triangle*] coordinates {(1000,287.3) (10000,250.4) (20000,494.5) (100000,499.3) (1000000,1401.5)};
  \addplot[crn, thick, mark=*] coordinates {(1000,328.3) (10000,292.5) (20000,418.6) (100000,695.1) (1000000,1612.8)};
  \addplot[cmap, thick, mark=square*] coordinates {(1000,225.8) (10000,343.6) (20000,766.8) (100000,689.0) (1000000,1822.6)};
  \addplot[cum, thick, mark=diamond*] coordinates {(1000,111.6) (10000,134.2) (20000,142.9) (100000,262.1) (1000000,708.8)};
  \legend{ABR, Rouge-noir, \texttt{map}, \texttt{unordered\_map}}
\end{loglogaxis}
\end{tikzpicture}
\caption{Suppression, clés en ordre aléatoire (échelles logarithmiques).}
\end{figure}

\section{Analyse}
\begin{itemize}
  \item \textbf{Hauteur} : avec des clés mélangées, l'ABR reste assez équilibré (50 niveaux pour un million de clés), mais le rouge-noir est \textbf{deux fois moins haut} (24 niveaux, proche de $\log_2 10^6 \approx 20$).
  \item \textbf{Recherche} : le rouge-noir est le plus rapide des arbres pour toutes les tailles. À $10^6$, il cherche \textbf{1,3 fois plus vite} que l'ABR (708 contre 896~ns) : la recherche ne fait aucune réparation et profite entièrement de la hauteur plus faible.
  \item \textbf{Insertion} : les trois arbres sont proches. \texttt{map} est devant pour les petites tailles ($10^3$, $10^4$), l'ABR à $2\times10^4$ et $10^5$, et notre rouge-noir à $10^6$ (1\,365 contre 1\,398~ns pour \texttt{map} et 1\,551~ns pour l'ABR).
  \item \textbf{Suppression} : l'ABR est le plus rapide des arbres dans la plupart des cas (1\,402 contre 1\,613~ns pour le rouge-noir à $10^6$, soit 13~\% de moins). Avec des clés mélangées, il reste assez bas, et il n'a \textbf{aucune réparation} à faire, alors que le rouge-noir doit recolorer et parfois tourner après avoir retiré un nœud noir.
  \item Pour le rouge-noir et \texttt{map}, la suppression est l'opération la plus chère à $10^6$ : il faut chercher le nœud, puis son successeur, puis réparer. Pour l'ABR, qui ne répare rien, elle coûte moins cher que l'insertion.
  \item Les mesures sous Windows sont un peu bruitées : par exemple, la recherche dans l'ABR à $2\times10^4$ (445~ns) est plus lente qu'à $10^5$ (304~ns), ce qui n'a pas de raison théorique.
\end{itemize}

% =====================================================================
\chapter{Résultats : clés en ordre croissant}

\section{Données brutes}
\begin{table}[H]
\centering
\caption{Ordre croissant : temps moyen par opération (ns). « -- » : non mesuré (trop long) ou sans objet.}
\small
\begin{tabular}{rlrrrr}
  \toprule
  $N$ & Structure & Hauteur & Insertion & Recherche & Suppression \\
  \midrule
  1\,000 & ABR & \cellcolor{ko} 1\,000 & \cellcolor{ko} 2\,850.5 & \cellcolor{ko} 3\,030.3 & \cellcolor{ko} 1\,205.1 \\
         & Rouge-noir & 17 & 355.0 & 53.8 & 233.8 \\
         & \texttt{map} & -- & 333.1 & 128.4 & 338.1 \\
         & \texttt{unordered\_map} & -- & 402.5 & 23.0 & 176.3 \\
  \midrule
  10\,000 & ABR & \cellcolor{ko} 10\,000 & \cellcolor{ko} 32\,272.7 & \cellcolor{ko} 38\,824.4 & \cellcolor{ko} 17\,775.7 \\
          & Rouge-noir & 24 & 240.7 & 116.7 & 263.1 \\
          & \texttt{map} & -- & 213.1 & 210.4 & 408.1 \\
          & \texttt{unordered\_map} & -- & 222.7 & 27.1 & 165.9 \\
  \midrule
  20\,000 & ABR & \cellcolor{ko} 20\,000 & \cellcolor{ko} 66\,420.1 & \cellcolor{ko} 84\,230.7 & \cellcolor{ko} 44\,145.6 \\
          & Rouge-noir & 26 & 272.3 & 184.2 & 502.7 \\
          & \texttt{map} & -- & 236.0 & 274.0 & 524.5 \\
          & \texttt{unordered\_map} & -- & 171.7 & 23.0 & 152.7 \\
  \midrule
  100\,000 & Rouge-noir & 31 & 321.1 & 323.0 & 572.2 \\
           & \texttt{map} & -- & 211.1 & 463.2 & 652.7 \\
           & \texttt{unordered\_map} & -- & 198.3 & 36.2 & 278.9 \\
  \midrule
  1\,000\,000 & Rouge-noir & 37 & 407.6 & 661.8 & 1\,418.5 \\
              & \texttt{map} & -- & 358.0 & 1\,314.4 & 1\,541.9 \\
              & \texttt{unordered\_map} & -- & 247.8 & 81.4 & 595.1 \\
  \bottomrule
\end{tabular}
\end{table}

\section{Visualisation}
\begin{figure}[H]
\centering
\begin{tikzpicture}
\begin{loglogaxis}[width=0.68\linewidth, height=6.5cm, xlabel={$N$}, ylabel={ns par insertion},
  legend style={font=\small, at={(1.03,0.5)}, anchor=west}, grid=major, grid style={black!12}]
  \addplot[cabr, thick, mark=triangle*] coordinates {(1000,2850.5) (10000,32272.7) (20000,66420.1)};
  \addplot[crn, thick, mark=*] coordinates {(1000,355.0) (10000,240.7) (20000,272.3) (100000,321.1) (1000000,407.6)};
  \addplot[cmap, thick, mark=square*] coordinates {(1000,333.1) (10000,213.1) (20000,236.0) (100000,211.1) (1000000,358.0)};
  \addplot[cum, thick, mark=diamond*] coordinates {(1000,402.5) (10000,222.7) (20000,171.7) (100000,198.3) (1000000,247.8)};
  \legend{ABR, Rouge-noir, \texttt{map}, \texttt{unordered\_map}}
\end{loglogaxis}
\end{tikzpicture}
\caption{Insertion, clés en ordre croissant : l'ABR dégénère, les autres restent plats.}
\end{figure}

\section{Analyse}
\begin{itemize}
  \item \textbf{L'ABR dégénère} : chaque clé est plus grande que toutes les précédentes et part toujours à droite. L'arbre devient une liste (hauteur $= N$), et quand $N$ double, le temps par opération double aussi. À $N = 20\,000$, il est \textbf{environ 240 fois plus lent} à l'insertion que le rouge-noir (66\,420 contre 272~ns), \textbf{460 fois} à la recherche (84\,231 contre 184~ns) et \textbf{90 fois} à la suppression (44\,146 contre 503~ns).
  \item La suppression de l'ABR souffre un peu moins : les clés sont supprimées dans un ordre aléatoire, et la liste raccourcit au fur et à mesure.
  \item \textbf{Le rouge-noir reste équilibré} : 37 niveaux pour un million de clés, sous la borne $2\log_2(n+1) \approx 40$. Les rotations corrigent la pente à mesure.
  \item \texttt{map} insère un peu plus vite que notre rouge-noir dans tous les cas en ordre croissant (358 contre 408~ns à $10^6$), mais notre rouge-noir reste plus rapide en recherche et en suppression.
  \item L'insertion croissante est \textbf{plus rapide} qu'en ordre aléatoire pour les arbres équilibrés (408 contre 1\,365~ns à $10^6$ pour le rouge-noir) : la descente suit toujours le même chemin, à droite, que le processeur garde en cache et prévoit sans erreur.
\end{itemize}

% =====================================================================
\chapter{Analyse comparative approfondie}

\section{Notre arbre rouge-noir face à \texttt{std::map}}
Les deux structures sont des arbres rouge-noir construits avec le même algorithme : elles ont la même forme et la même hauteur. Pourtant, notre version \textbf{cherche environ 2 fois plus vite} (708 contre 1\,386~ns à $10^6$ en aléatoire, 662 contre 1\,314~ns en croissant), et, à partir de $2\times10^4$ clés, elle supprime aussi plus vite (1\,613 contre 1\,823~ns à $10^6$). \texttt{map} n'est devant qu'à l'insertion (petites tailles et ordre croissant) et à la suppression de $1\,000$ clés.

L'écart de recherche vient du code machine produit par le compilateur :
\begin{itemize}
  \item Notre boucle \texttt{c = (v < c->valeur) ? c->gauche : c->droite;} est traduite en une instruction \texttt{cmov} (déplacement conditionnel) : le choix entre gauche et droite se fait \textbf{sans branchement}, et le processeur n'a rien à deviner.
  \item \texttt{map::find} descend avec un \texttt{if} que le processeur doit prédire. Avec des clés aléatoires, il se trompe environ une fois sur deux, et chaque erreur coûte une quinzaine de cycles.
  \item \texttt{map::find} descend toujours jusqu'en bas de l'arbre (recherche de type \texttt{lower\_bound}), alors que notre recherche s'arrête dès qu'elle trouve la clé.
  \item Un nœud de \texttt{map<int,int>} est plus gros : il stocke une paire clé-valeur.
\end{itemize}
\texttt{std::map} reste plus générale : elle accepte n'importe quel type et n'importe quel comparateur, et fournit des itérateurs. Notre arbre ne stocke que des \texttt{int}.

\section{ABR ou rouge-noir : le prix de l'équilibrage}
En ordre aléatoire, l'équilibrage n'apporte qu'un gain modeste : le rouge-noir cherche 1,3 fois plus vite, insère à peu près aussi vite, et \textbf{supprime même plus lentement} que l'ABR, à cause de la réparation. En ordre croissant, en revanche, il évite un effondrement d'un facteur 90 à 460. L'équilibrage est donc une \textbf{assurance} : un petit coût dans le cas favorable, contre une protection totale dans le pire cas.

\section{Les arbres face à la table de hachage}
\texttt{unordered\_map} est la plus rapide presque partout : à $10^6$ en aléatoire, elle cherche environ \textbf{5,5 fois plus vite} que le rouge-noir (129 contre 708~ns), insère et supprime environ 2 fois plus vite. Elle calcule directement la case de la clé au lieu de descendre une vingtaine de niveaux. Seule exception : avec $1\,000$ clés en aléatoire, son insertion (419~ns) n'est pas meilleure que celle des arbres, car le coût de création et d'agrandissement de la table pèse plus lourd que le gain du hachage.

Mais elle ne garde \textbf{aucun ordre} : impossible d'obtenir les clés triées, le minimum, le maximum ou toutes les clés d'un intervalle. Son temps $O(1)$ n'est qu'une moyenne : avec beaucoup de collisions, il devient $O(n)$, et l'agrandissement de la table recopie tout d'un coup.

\section{Mémoire}
\begin{table}[H]
\centering
\caption{Taille d'un nœud sur une machine 64 bits (hors surcoût de \texttt{new}).}
\small
\begin{tabular}{lll}
  \toprule
  Structure & Contenu d'un nœud & Taille \\
  \midrule
  ABR & clé + 2 pointeurs & 24 octets \\
  Rouge-noir & clé + couleur + 3 pointeurs & 32 octets \\
  \texttt{map<int,int>} & couleur + 3 pointeurs + paire clé-valeur & 40 octets \\
  \texttt{unordered\_map<int,int>} & pointeur suivant + paire, + un tableau de cases & 16 octets + cases \\
  \bottomrule
\end{tabular}
\end{table}
L'équilibrage coûte un pointeur \texttt{parent} et une couleur par nœud, soit 8~Mo de plus pour un million de clés : un prix faible face au gain de hauteur.

\section{Bilan}
\begin{table}[H]
\centering
\caption{Quelle structure choisir ?}
\small
\begin{tabular}{p{0.30\linewidth}p{0.28\linewidth}p{0.32\linewidth}}
  \toprule
  Besoin & Structure & Raison \\
  \midrule
  Savoir si une clé existe & \texttt{unordered\_map} / \texttt{unordered\_set} & $O(1)$ en moyenne, la plus rapide \\
  Clés triées, minimum, intervalles & \texttt{map} / \texttt{set} & arbre rouge-noir, $O(\log n)$ garanti \\
  Données pouvant arriver triées & jamais un ABR simple & dégénère en liste \\
  Comprendre l'équilibrage & arbre écrit à la main & tout le fonctionnement est visible \\
  \bottomrule
\end{tabular}
\end{table}

% =====================================================================
\chapter{Conclusion}

\section{Synthèse des résultats}
\begin{enumerate}
  \item Notre arbre rouge-noir est correct : il donne les mêmes arbres que les exercices faits à la main, et ses règles sont vérifiées automatiquement après les insertions et après les suppressions.
  \item Avec des clés triées, l'ABR devient une liste : à $N = 20\,000$, il est environ \textbf{240 fois plus lent} à l'insertion, \textbf{460 fois} à la recherche et \textbf{90 fois} à la suppression que le rouge-noir, qui garde une hauteur de 37 pour un million de clés.
  \item Avec des clés mélangées, le rouge-noir \textbf{cherche 1,3 fois plus vite} que l'ABR, insère à peu près aussi vite, mais supprime environ 13~\% plus lentement, à cause de la réparation des couleurs.
  \item Notre rouge-noir cherche environ \textbf{2 fois plus vite que \texttt{std::map}}, grâce à une boucle de recherche sans branchement et à un arrêt dès que la clé est trouvée ; \texttt{map} est un peu plus rapide à l'insertion en ordre croissant.
  \item \texttt{unordered\_map} est la plus rapide (environ 5,5 fois plus rapide que le rouge-noir pour chercher), mais ne garde aucun ordre.
\end{enumerate}

\section{Recommandations}
\begin{itemize}
  \item Pour un dictionnaire sans besoin d'ordre : \texttt{unordered\_map}.
  \item Pour des données triées ou des requêtes par intervalle : \texttt{map}, qui garantit $O(\log n)$ quel que soit l'ordre d'arrivée.
  \item Un ABR simple ne convient que si l'on sait que les données arrivent dans un ordre aléatoire.
  \item Pour des mesures plus stables, répéter le benchmark davantage de fois et fermer les autres programmes : les mesures sous Windows sont un peu bruitées.
\end{itemize}

\section{Conclusion finale}
L'équilibrage ne rend pas chaque opération plus rapide : en ordre aléatoire, l'ABR insère aussi vite que le rouge-noir et supprime même plus vite. Son vrai rôle est de donner une \textbf{garantie} : quel que soit l'ordre d'arrivée des données, la hauteur reste en $O(\log n)$. Cette garantie évite l'effondrement de l'ABR face à des données triées, pour un coût de quelques recolorations, d'au plus deux rotations par insertion et trois par suppression, et d'un pointeur de plus par nœud.

\end{document}
